\section{Production Deployment Considerations}

Protocol selection for production Byzantine fault tolerant systems is rarely settled by a single figure of merit. The properties that bind in practice depend on the deployment envelope: the number of validators the operator expects to admit, whether membership is closed or open, the latency budget imposed by the application, and the adversarial model the system must survive. This section maps the structural characteristics established above onto those axes and identifies, for each deployment class, the protocol feature that is actually decisive.

\subsection{Validator Count and Communication Complexity}

The dominant scaling constraint is communication complexity. Classical PBFT requires $\mathcal{O}(n^2)$ messages per consensus instance~\cite{castro1999practical}, a cost that is effectively inert for small validator sets and prohibitive as membership grows. Deployments should therefore treat validator count as a design parameter rather than an emergent property: the quadratic term determines the point at which message load, not execution or storage, becomes the throughput ceiling. Orthogonal to complexity is the resilience bound itself. Because safety requires $n > 3f$, each additional tolerated Byzantine fault costs three additional validators, and the region $n \leq 3f$ is unavailable to any deterministic protocol irrespective of timing assumptions or communication pattern. Operators consequently face a fixed exchange rate between fault tolerance and coordination cost. Protocols that aggregate votes into quorum certificates and drive a linear per-view message pattern, as in HotStuff~\cite{yin2019hotstuff}, relax the quadratic term by concentrating aggregation work at the leader, which shifts the bottleneck from network fan-out to leader bandwidth and signature verification.

\subsection{Permissioned versus Permissionless Membership}

Permissioned deployments enjoy two properties that materially simplify protocol choice: validator identities are known in advance, and network behaviour can be measured. Timeouts can therefore be calibrated against observed delay distributions rather than assumed conservatively, making the partially synchronous guarantee of PBFT-lineage protocols operationally tight. Permissionless deployments forfeit both properties. Membership churns, delay bounds cannot be calibrated against a stable population, and the relevant safety question shifts from whether agreement holds to whether violations can be attributed and penalised. Casper~\cite{buterin2020combining} addresses this by composing a finality gadget over a separate proposal mechanism, decoupling chain growth from finalisation so that open participation in block production coexists with a bounded finalising committee. Tendermint~\cite{buchman2016tendermint} occupies the intermediate position: validator sets are open by stake yet deliberately bounded in cardinality, which preserves a PBFT-style commit structure while admitting rotating membership.

\subsection{Latency Sensitivity and Finality Semantics}

Partially synchronous protocols provide deterministic termination only after the Global Stabilization Time; before it, no latency bound holds. Latency-sensitive deployments must therefore budget against view change cost rather than steady-state throughput. Aggressive timeouts shorten recovery when a leader fails but provoke spurious view changes that waste rounds; conservative timeouts protect the common case at the expense of tail latency during genuine faults. Where leaders rotate every view, the per-view cost of leader replacement is paid continuously rather than exceptionally, so protocols whose view change retains quadratic cost amortise poorly under adversarial leader scheduling. DAG-based designs invert the tradeoff by separating dissemination from ordering: throughput tracks available bandwidth instead of a single leader's critical path, at the cost of additional ordering latency.

\subsection{Decisive Features by Deployment Class}

For small closed validator sets under hard latency ceilings, such as replicated databases and matching engines, the decisive feature is immediate deterministic finality with calibrated timeouts; the quadratic message term is not binding at that scale. For moderate open-stake validator sets, the decisive feature is bounded per-view authenticator cost together with pipelined certificate formation. For large adversarial memberships, the decisive feature is accountable safety and economic finality rather than raw throughput. For throughput-dominated workloads tolerant of ordering delay, bandwidth-proportional dissemination dominates.

No single protocol dominates across these classes. The engineering discipline lies in identifying which constraint binds first and selecting the variant whose asymptotics fail last along that axis.
